\section{感悟}
通过第二周的实验练习，形成以下认识和体会：
\begin{itemize}
\item 任务控制与信号的可靠性：通过 trap 与 SIGTERM 的配合，可以让后台任务在被终止时执行清理逻辑；jobs 提供的工作号查询也避免了手工抄写 PID 带来的错误。
\item 语言服务器对语义重构的价值：借助编辑器的跳转到定义、查找引用和重命名符号，一次重命名即可让定义与所有引用同步更新，显著降低了手动替换的遗漏风险。
\item 调试器带来的定位效率：相较于反复插入 print，断点、单步与变量观察能更直接地揭示缺陷根因；本次归并排序的问题只是把 right[j] 误写成了 right[i]。
\item 先测量、再优化的方法论：cProfile 明确指出瓶颈在于 list 的线性 in 查找；将去重从 list 改为 set 后，复杂度由 O($n^2$) 降为 O(n)，在 30000 词规模下获得约 7 倍加速。
\end{itemize}

命令行环境负责进程与信号管理，开发工具与语言服务器提升编码与重构效率，调试器与性能分析器则分别从正确性和性能两个维度保障代码质量。三者共同构成了从“能运行”走向“正确、高效、可维护”的完整工具链。

\subsection*{Github提交}
此处上传Github提交记录。(报告于完成后上传)
\begin{figure}[H]
    \centering
    \includegraphics[width=0.8\linewidth]{git}
\end{figure}
